% STATUS: written 2026-11-06 (English gate clean; sources read in full as recorded below).
% Sources read in full for the cited passages: van den Essen-Kuroda-Koenig 2021 (pp. 91, 101, 104, 129--131), Hauser-Mueller 1998 (p. 211).
\section{The Jacobian conjecture}\label{sec:commutative:jacobian}

Section~\ref{sec:commutative:locally-nilpotent} ended with open problems of
polynomial rings. The most famous of them is the Jacobian
conjecture~\cite[p.~91]{VanEssenKurodaKoenig2021}. The story of derivations
meets it directly~\cite[p.~91]{VanEssenKurodaKoenig2021}.

Call a polynomial map $F\colon k^n \to k^n$ a \emph{Keller map} when
$\det JF = 1$ \cite[p.~106]{VanEssenKurodaKoenig2021}. The conjecture asks
whether every Keller map is invertible
\cite[p.~106]{VanEssenKurodaKoenig2021}.

The conjecture has many equivalent forms, several of which speak the language
of this book \cite[p.~91]{VanEssenKurodaKoenig2021}. Van den Essen, Kuroda and
Koenig open their chapter on the new results with the words ``there are
various equivalent formulations of the Jacobian Conjecture''
\cite[p.~91]{VanEssenKurodaKoenig2021}.

This section recalls three of them that involve
derivations~\cite[p.~91]{VanEssenKurodaKoenig2021}. Each one turns the
statement about polynomial maps into a statement about one derivation, or one
Lie algebra of derivations~\cite[p.~91]{VanEssenKurodaKoenig2021}.

The first lives on the plane
\cite[Example~5.2.14, pp.~130--131]{VanEssenKurodaKoenig2021}. The divergence of a derivation is the object of study
here~\cite[p.~130]{VanEssenKurodaKoenig2021}. For a $k$-derivation
$D = a\,\partial_x + b\,\partial_y$ of $k[x, y]$, the \emph{divergence} of $D$
is $\operatorname{div} D = \partial_x a + \partial_y b$
\cite[p.~130]{VanEssenKurodaKoenig2021}.

A \emph{Mathieu-Zhao space}, or MZ-space, is a subspace whose radical and
strong radical coincide \cite[Definition~5.2.1, p.~123]{VanEssenKurodaKoenig2021}.
Here the reader only needs one property from the proof of
Example~\ref{prop:comm:jc-div}, stated next. An MZ-space that contains $1$ is
the whole ring \cite[Example~5.2.11, pp.~129--130]{VanEssenKurodaKoenig2021}.

\begin{proposition}[{\cite[Example~5.2.14, pp.~130--131]{VanEssenKurodaKoenig2021}}]\label{prop:comm:jc-div}
The following two conditions are equivalent.
\begin{enumerate}
\item For every $k$-derivation $D$ of $k[x,y]$ such that $1$ lies in the
image of $D$ and $\operatorname{div} D = 0$, the image of $D$ is an
MZ-space of $k[x,y]$.
\item The two-dimensional Jacobian conjecture is true.
\end{enumerate}
\end{proposition}

The proof recalled in the source builds a derivation from the map itself
\cite[pp.~130--131]{VanEssenKurodaKoenig2021}. For $F = (f, g)$ with
$\det JF = 1$ it sets
\[
  D = g_y\,\partial_x - g_x\,\partial_y .
\]
It notes that $\operatorname{div} D = 0$ and $1 = \det JF = Df$
\cite[pp.~130--131]{VanEssenKurodaKoenig2021}.

It then shows that $D$ is locally nilpotent with kernel $k[g]$
\cite[pp.~130--131]{VanEssenKurodaKoenig2021}. So a Keller map in two
variables is invertible exactly when the image of this derivation is a Mathieu
space of the ring~\cite[pp.~130--131]{VanEssenKurodaKoenig2021}.

The second formulation uses a Lie algebra of the kind developed in
Section~\ref{sec:commutative:vector-fields-geometry}
\cite[p.~211]{HauserMuller1998}. Let $K$ be an algebraically closed field of
characteristic zero. A polynomial vector field has the form
\[
  \sum_i a_i\,\partial_{x_i}, \qquad a_i \in K[x_1,\dots,x_n].
\]
Let $\mathcal{P}$ be the Lie algebra of these fields
\cite[p.~211]{HauserMuller1998}. Let $\mathcal{D}$ be the Lie algebra of formal
vector fields, whose coefficients are formal power
series~\cite[p.~211]{HauserMuller1998}.

\begin{proposition}[{\cite[Proposition~1, p.~211]{HauserMuller1998}}]\label{prop:comm:hm}
The following statements are equivalent.
\begin{enumerate}
\item Every endomorphism $\varphi$ of $K[x_1,\dots,x_n]$ with invertible
Jacobian matrix is an automorphism.
\item Every injective endomorphism of $\mathcal{P}$ is an automorphism.
\item Every automorphism $\Phi$ of $\mathcal{D}$ with
$\Phi(\mathcal{P}) \subset \mathcal{P}$ satisfies
$\Phi(\mathcal{P}) = \mathcal{P}$.
\end{enumerate}
\end{proposition}

Hauser and Müller credit the equivalence of the first two statements to
Kulikov~\cite[p.~211]{HauserMuller1998}. Thus the Jacobian conjecture becomes a
property of one infinite-dimensional Lie algebra of
derivations~\cite[Proposition~1, p.~211]{HauserMuller1998}.

The polynomial vector fields must be
preserved~\cite[Proposition~1, p.~211]{HauserMuller1998}. They are preserved
under the automorphisms of the formal ones that map them into
themselves~\cite[Proposition~1, p.~211]{HauserMuller1998}.

The third formulation links the conjecture to two other great open
problems~\cite[Theorem~4.4.1, p.~101]{VanEssenKurodaKoenig2021}. The Dixmier
conjecture asks whether every endomorphism of a Weyl algebra is an
automorphism~\cite[Theorem~4.4.1, p.~101]{VanEssenKurodaKoenig2021}. The
Poisson conjecture makes the same demand of the Poisson
algebra~\cite[Theorem~4.4.1, p.~101]{VanEssenKurodaKoenig2021}.

\begin{theorem}[{\cite[Theorem~4.4.1, p.~101]{VanEssenKurodaKoenig2021}}]\label{thm:comm:dc-jc-p}
For each $n \ge 1$ one has the implications
$\mathrm{DC}(n,\mathbb{C}) \Rightarrow \mathrm{JC}(n,\mathbb{C})$,
$\mathrm{JC}(2n,\mathbb{C}) \Rightarrow \mathrm{P}(n,\mathbb{C})$ and
$\mathrm{P}(n,\mathbb{C}) \Rightarrow \mathrm{DC}(n,\mathbb{C})$; that is,
the Dixmier, Jacobian, and Poisson conjectures are equivalent.
\end{theorem}

The implication from the Dixmier conjecture to the Jacobian conjecture was
classical \cite[p.~91]{VanEssenKurodaKoenig2021}. Tsuchimoto proved the
converse in 2005, and Belov-Kanel and Kontsevich gave an independent proof
\cite[p.~91]{VanEssenKurodaKoenig2021}.

The conjecture of this section is open, and this section states it in the form
the book uses. % claim-ok: own signpost introducing the moved problem, which carries its own citation

\begin{problem}[{\cite[p.~81]{Freudenburg2017}}]\label{prob:open:jc}
Let $F\colon k^n\to k^n$ be a polynomial map over a field $k$ of characteristic
$0$ whose Jacobian determinant is a nonzero constant.
Is $F$ then an automorphism?
\end{problem}

Freudenburg records the problem in the language of derivations
\cite[p.~81]{Freudenburg2017}. A jacobian derivation of $k[x_1,\dots,x_{n-1}]$
has a slice exactly when it is locally nilpotent and its kernel is a polynomial
ring in $n-1$ variables \cite[p.~81]{Freudenburg2017}. The problem is open for
every $n\ge2$ \cite[pp.~1--2]{BelovKanelKontsevich2007}.

Theorem~\ref{thm:comm:dc-jc-p} also ties the Jacobian story to the Weyl
algebra, and Chapter~12 of this book develops that side.
Section~\ref{sec:polynomial-weyl:weyl} states the Dixmier problem.
Section~\ref{sec:commutative:vector-fields-geometry} follows the derivations
of $k[x_1,\dots,x_n]$ to their smooth limit, the vector fields on a manifold.
